% !TEX program = xelatex
\documentclass[11pt,letterpaper]{article}

% === GEOMETRY === (narrow measure, essay style)
\usepackage[top=1in, bottom=1in, left=1.15in, right=1.15in]{geometry}

% === FONTS ===
\usepackage{fontspec}
\usepackage{unicode-math}

\setmainfont{Latin Modern Roman}[
  Extension=.otf,
  UprightFont=lmroman10-regular,
  BoldFont=lmroman10-bold,
  ItalicFont=lmroman10-italic,
  BoldItalicFont=lmroman10-bolditalic,
]
\setsansfont{Latin Modern Sans}[
  Extension=.otf,
  UprightFont=lmsans10-regular,
  BoldFont=lmsans10-bold,
  ItalicFont=lmsans10-oblique,
  BoldItalicFont=lmsans10-boldoblique,
]
\setmonofont{Latin Modern Mono}[
  Extension=.otf,
  UprightFont=lmmono10-regular,
  ItalicFont=lmmono10-italic,
  Scale=0.85,
]

% === PACKAGES ===
\usepackage{xcolor}
\usepackage{titlesec}
\usepackage{enumitem}
\usepackage{fancyhdr}
\usepackage{hyperref}
\usepackage{xurl}
\usepackage{ragged2e}
\usepackage{microtype}

\usepackage{../../papers/essay-may-26/ac-paper-essay}

\hypersetup{
  colorlinks=true,
  linkcolor=acpurple,
  urlcolor=acpurple,
  pdfauthor={Jeffrey Scudder / @jeffrey},
  pdftitle={Scores for Social Software --- Session \#5 Brief},
}

\pagestyle{fancy}
\fancyhf{}
\renewcommand{\headrulewidth}{0pt}
\fancyhead[C]{\footnotesize\color{acgray}\thepage}

\newcommand{\ac}{Aesthetic\textcolor{acpink}{.}Computer}
\newcommand{\acos}{AC Native OS}
\newcommand{\np}{notepat}

\begin{document}

\thispagestyle{empty}
\vspace*{0.4in}

\essaytitle{Session \#5 Brief}
\essaydate{Scores for Social Software, Cycle 2 \enspace$\cdot$\enspace 28 May 2026 \enspace$\cdot$\enspace Jeffrey Scudder / @jeffrey}

In March I submitted \ac{} not as a thing I would build, but as a thing I was already building. The score was to develop the platform in the open across the cycle and let the cohort engage with the live codebase, the live community, and the design choices in motion. Two months in, the score has held. The platform did not stop to perform for the cycle. The cycle joined the platform as it was already moving.

What follows is a brief written for the conversation Casey called for tomorrow: a status check on the individual score, the publication, and the performance. I have tried to keep it modal --- what \emph{has} happened is reported, what \emph{may} happen at Fuser on 13 June is sketched honestly rather than promised.

\section{Where the Individual Score Stands}

The score, restated in one line: \emph{develop \ac{} in the open over ten weeks, share the decisions, and let the cohort participate as users and critics.}

The open part is observable. Since Session \#4 (14 May) alone, the public repository has taken 202 commits. The cohort has had access to every change, every dossier, every paper, every release as it landed. The development was not staged for the cycle; it was the cycle, in plain view.

The concrete deltas, since the March proposal:

\begin{itemize}
  \item \textbf{The platform itself.} \ac{} continues to operate. The network grew from 2,798 to roughly 2,812 @handles; KidLisp programs from 16,174 to 16,779; pieces, paintings, chats and moods all continued to accumulate at the cadence the proposal described.
  \item \textbf{New lanes that did not exist in March.} The \texttt{pop/} music lane shipped its first single (\textit{trancenwaltz}, 17 May, DistroKid) and now carries multiple parallel tracks --- chillwave, trancepenta, marimba, hellsine, helpabeach, gradus, hippyhayzard --- each with its own synth language, vocal pipeline, and visualizer.
  \item \textbf{\acos{}} now boots on surplus ThinkPads in 7.3 seconds. The OTA build pipeline runs remotely on a dedicated server (\texttt{oven}). The hardware story is real.
  \item \textbf{The papers mill} grew from a handful of arxiv-style entries into a structured platter: thirty-plus arxiv papers, seventeen institutional dossiers, an essay lane, and a card-format printable shadow of every paper.
  \item \textbf{Infrastructure adjacent to the social network.} Slab (the macOS menubar surface), Menuband (the floating KidLisp piano), Wave Wizard, the agent-memory pipeline, the \texttt{oven} render service --- all built or extended in the open.
  \item \textbf{External engagements.} CultureHub LA residency submitted; Tezos Foundation ecosystem grant submitted (KidLisp Keeps v12); CAC IAF Legacy logged; UCLA Social Software invoice closed.
\end{itemize}

What I would offer the room is not a triumphant list. It is the texture of \emph{developing in public as the score itself.} The honest question is whether the cohort experienced this as participatory, or whether the firehose was too wide to engage with critically. I want feedback on that specifically.

\section{The Publication}

The proposal anticipated that the platform's existing publishing infrastructure --- the papers mill, the cards, the dossiers --- would carry the cycle's documentary load. That has been true.

My contribution to the cohort publication, as I'd propose it for the wrap-up:

\begin{itemize}
  \item \textbf{A printed score} for \ac{} as an open-development project: a single-sheet card stating the score in the form the platform already uses (the cards lane). The card is reproducible from \texttt{.tex} source; the cohort gets the file, not just the print.
  \item \textbf{The platform statistics snapshot} from the date of the publication, in the same tabular form as the proposal's Appendix A. The snapshot is the documentary evidence of what \emph{the cycle's room of users} produced inside the platform during those ten weeks --- not a marketing figure, an observational one.
  \item \textbf{One short essay} in the magazine-essay register I have been developing in \texttt{papers/essay-*/}, written for the publication: an honest reflection on what it meant to develop the platform alongside a cohort that was not the platform's typical audience. Companion to \textit{Aesthetic May '26}, but tighter.
  \item \textbf{Optional:} A KidLisp card --- a single line of code that paints itself, printed on the back of the score. The simplest legible artifact the system produces.
\end{itemize}

What I'd like the room to weigh in on: whether the publication wants a unified design across contributors, or whether each contributor brings their own object and the publication is an assemblage. I have a strong design grammar (\texttt{ac-paper-essay.sty}, the AC palette, YWFT Processing display + Latin Modern body) and I am happy to either submit inside someone else's frame or offer mine as one option.

\section{The Performance / Presentation}

The final event is at Fuser in Frogtown on 13 June, noon--2pm. I have not yet been to the room. What I can offer, in rough order of how performable I believe each is on that day:

\begin{itemize}
  \item \textbf{\np{} live.} \np{} is the platform's front door and, by line count, its largest single piece. It has polyphony, room reverb, octave shift, drum routing, and a menubar visualization. A five-to-eight-minute set on \np{} is the most rehearsed thing I could bring. I would play, narrate the instrument as I played, and end with a piece written that morning in front of the room.
  \item \textbf{An \acos{} boot.} A surplus ThinkPad, no operating system on the disk, USB plugged in, power on. Seven seconds later the room watches \ac{} appear on bare metal. Small, but the point of the project lives inside those seven seconds.
  \item \textbf{A KidLisp set.} I write KidLisp live, on screen, and let the room call out modifications. The same way one might run a small Scheme demonstration, but for an audience that is not necessarily programming-literate. The cards from the publication double as a printed program for this section.
  \item \textbf{A \texttt{pop/} cut with the visualizer.} The strongest argument that the platform is more than its IDE. The risk is that it shifts the register from \emph{social software} toward \emph{audio-visual performance}; the question is whether that shift serves the cohort or distracts.
  \item \textbf{An open prompt.} I sit at the laptop, the prompt is on the screen, and the room types into it. Whatever the room types, the platform runs. The most truthful demonstration of what \ac{} is. Also the hardest to bound in a fixed window.
\end{itemize}

The most honest read is that I can deliver any of these well, and I would rather pick one with the room than guess. My current preference, if the room has no opinion: \np{} as the spine, with the \acos{} boot as the opening gesture and a KidLisp card on every seat.

\section{Open Questions for Tomorrow}

I want to leave time for these.

\begin{itemize}
  \item Did the open-development format land as participatory, or was the throughput of the project hard to read into as a cohort?
  \item Does the cohort publication want unified design or assemblage?
  \item For Fuser on 13 June: one performer at a time, or some interleaved or simultaneous format? Is there a shared piece we want to attempt as a group?
  \item Is there a part of the platform any of you wanted to engage with and didn't have a hook into? I can build the hook this week if there is one.
\end{itemize}

\section{A Note on What I'm Not Bringing}

The proposal said this would be the cycle where I composed in conversation rather than alone. The honest report is that I have continued to compose at the rate I always compose at, and the cohort's voices entered the work alongside many others. That may have been the right calibration, or it may have meant the cycle's signal was harder to find. I want to know which.

The two weeks to wrap-up are enough to refactor any of this. I'd rather hear your read tomorrow and rebuild the contribution to match, than ship what I already have.

\colophon{Prepared for Session \#5 \enspace$\cdot$\enspace Zoom, noon--2pm \enspace$\cdot$\enspace aesthetic.computer \enspace$\cdot$\enspace ORCID: \href{https://orcid.org/0009-0007-4460-4913}{0009-0007-4460-4913}.}

\end{document}
